% !TEX root = ../../main.tex
%% Défi du chapitre 16 — Les nombres de Stirling de seconde espèce
%% Résultats vérifiés par verifications/defis_16_20.py.

\rubrique{Défi}

\begin{defi}{Les nombres de Stirling de seconde espèce}{1 h 30}
Pour $n\geqslant1$ et $k\geqslant1$, on note $S(n,k)$ le nombre de façons de
partager l'ensemble $\intFF{1}{n}$ en $k$ parties non vides (des
\emph{blocs}), sans tenir compte de l'ordre des blocs ; on convient que
$S(0,0)=1$, et $S(n,k)=0$ si $k>n$ ou si $k=0<n$. Ainsi $S(3,2)=3$ :
$\ens{1}\ens{2,3}$, $\ens{2}\ens{1,3}$, $\ens{3}\ens{1,2}$.

\partieExo{A}{Premières valeurs et récurrence}
\begin{questions}
  \item Déterminer $S(n,1)$, $S(n,n)$, puis montrer que
    $S(n,n-1)=\dbinom n2$ et $S(n,2)=2^{n-1}-1$.
  \item En distinguant si l'élément $n$ forme un bloc à lui seul ou non,
    établir la relation
    \[
    S(n,k)=S(n-1,k-1)+k\,S(n-1,k)\qquad(n,k\geqslant1).
    \]
    Retrouver ainsi les lignes $n=5$ et $n=6$ du triangle de la figure.
  \item Dessiner les $S(4,2)$ partitions de $\intFF{1}{4}$ en deux blocs.
\end{questions}

\medskip
\begin{center}
\begin{tikzpicture}[x=1.15cm,y=0.62cm]
  \foreach \row [count=\n] in {{1},{1,1},{1,3,1},{1,7,6,1},{1,15,25,10,1},{1,31,90,65,15,1}}{%
    \node[font=\scriptsize,text=otMuted,anchor=east] at (-0.9,-\n) {$n=\n$};
    \foreach \v [count=\k] in \row {%
      \node[font=\small,text=otInk,inner sep=1.5pt] (s\n\k) at ({\k-1-(\n-1)/2+2.5},-\n) {$\v$};}}
  \draw[otGreenDark,line width=0.8pt,rounded corners=2pt] ($(s53)+(-0.28,-0.32)$) rectangle ($(s53)+(0.28,0.32)$);
  \draw[otNavy,line width=0.8pt,-{Stealth[length=1.6mm]}] (s42) -- node[left,font=\tiny,text=otNavy]{$\times1$} (s53);
  \draw[otNavy,line width=0.8pt,-{Stealth[length=1.6mm]}] (s43) -- node[right,font=\tiny,text=otNavy]{$\times3$} (s53);
\end{tikzpicture}\\[0.5mm]
\legende{Le triangle des $S(n,k)$ ($k$ croissant de gauche à droite) :
$S(5,3)=S(4,2)+3\,S(4,3)=7+3\times6=25$.}
\end{center}

\partieExo{B}{Surjections et factorielles descendantes}
\begin{questions}[resume]
  \item Montrer que le nombre de surjections de $\intFF{1}{n}$ sur
    $\intFF{1}{k}$ vaut $k!\,S(n,k)$.
  \item En classant les applications de $\intFF{1}{n}$ dans $\intFF{1}{x}$
    ($x\in\N$) selon leur image, montrer que
    \[
    x^{n}=\sum_{j=0}^{n}S(n,j)\,x(x-1)\cdots(x-j+1) .
    \]
  \item En déduire l'égalité de polynômes
    $X^{n}=\sum_{j=0}^{n}S(n,j)\,X(X-1)\cdots(X-j+1)$, et que les
    coordonnées de $X^{n}$ dans la base binomiale $(H_j)$ du défi du
    chapitre 12 sont les $j!\,S(n,j)$. Vérifier sur $X^{3}$.
\end{questions}

\partieExo{C}{Nombres de Bell et formule explicite}
\begin{questions}[resume]
  \item On note $B_n=\sum_kS(n,k)$ le nombre total de partitions de
    $\intFF{1}{n}$ ($B_0=1$). En considérant le bloc qui contient $n+1$,
    montrer que $B_{n+1}=\sum_{k=0}^{n}\binom nkB_k$, et calculer
    $B_1,\dots,B_6$.
  \item À l'aide de l'exercice 15 (surjections sur trois éléments), montrer
    que $S(n,3)=\dfrac{3^{n}-3\cdot2^{n}+3}{6}$ et le vérifier sur le
    triangle.
\end{questions}
\end{defi}
